\section{Discussion}
\label{sec:discussion}

\para{The direction of the intuition is right; the mechanism is wrong.} Audience modeling clearly
helps: both \cot and \simtom beat \plain, and \simtom---an explicit one-shot audience model---is
the best condition on \emph{every} metric, including the adaptation-gap control that specifically
measures genuine audience use. So the premise that LLMs write generically partly because they
model the reader weakly is supported. What fails is the specific proposed fix. Simulating the
reader every sentence does not just add nothing; it actively backfires, dragging the writing below
even plain generation on overall quality.

\para{Why does per-sentence simulation hurt?} Three mechanisms are visible in the outputs.
\emph{First, local reactivity destroys global craft.} Choosing each next sentence from the
reader's momentary reaction yields prose that is locally coherent but loses arc, voice, and payoff.
In one children's-story task, \simtom builds a whimsical arc with a delightful closing button
(``Death's Help Desk: Making Life Interesting Since Forever''), whereas the per-sentence version
stays plot-mechanical, keeps an adult phrase (``living forever sucks''), and drifts into heavier
themes (``watching everyone I loved grow old and leave'')---each next sentence locally reasonable,
the whole globally off-key for a child. \emph{Second, it adapts less, not more}
(\secref{sec:res-adapt}). A single rich audience model front-loads a coherent stance---``this is a
scared 8-year-old: stay warm, concrete, gentle''---that governs the entire piece, whereas
per-sentence reactions are noisier and pull the text in inconsistent directions, washing out
adaptation. \emph{Third, bloat.} The intervention runs about $15\%$ longer and spends the most
thinking tokens, yet yields worse writing---the opposite of the token-for-quality trade the
hypothesis assumed.

\para{Granularity is the key axis.} Our result directly echoes and extends
\citet{chakrabarty2025aislop}'s finding that more reasoning does not improve writing: even
\emph{audience-specific} reasoning fails to help when applied at fine (per-sentence) granularity,
while a coarse, one-shot audience model does help. The sweet spot is ``model the reader once,
richly,'' not ``re-model every sentence.'' This isolates granularity---rather than the presence or
absence of audience modeling---as the design variable that decides whether audience reasoning
helps or fragments the prose.

\para{Validity.} Several safeguards make the negative result credible. The judges are in different
families from the generators, so self-preference cannot explain the ranking; pairwise judging is
blind and order-balanced against position bias; only clean final text is judged, after a two-step
protocol eliminated a real scaffolding-leak bug we caught mid-study (\secref{sec:method}); the
adaptation-gap control rules out a generic-quality confound; and the full ordering replicates
across two generators. The effect sizes are large (win-rates as low as $0.083$) and consistent,
not marginal.

\para{Limitations.} We test a single granularity---``every sentence.'' An intermediate granularity
(per-paragraph, or once-then-one-revision-pass) might recover a benefit; we did not sweep it. Fit
and quality come from LLM judges rather than humans; although the judges are cross-family and the
effect is large and replicated, a small human evaluation would strengthen the quality claim. The
result also bounds \emph{one} careful, faithful phrasing of the per-sentence prompt, not every
conceivable one. We study short-form writing ($\sim$150 words), where global consistency is easy
to maintain without per-sentence tracking; the intervention might matter more for long documents.
Finally, our task set is 60 tasks (48 for replication); paired tests and bootstrap intervals
mitigate this, but larger, more diverse task sets would tighten the estimates.

\para{Broader implications.} For practitioners, the actionable takeaway is simple and cheap:
prompt the model to build one rich picture of the reader before writing, and do not ask it to
re-simulate the reader sentence by sentence. More broadly, the result is a caution against
importing scaffolds that help \emph{comprehension} (where per-sentence perspective tracking is
known to help~\citep{perceptom2024}) directly into \emph{generation}, where global structure is
part of the quality being measured.
